\section{Smallest Multiple --- Problem 5}
\label{prob:5}
\noindent\textbf{Source:} \url{https://github.com/KO0I/projectEuler/blob/reorganize/solutions-by-language/solutions/c/smallest_multiple_5.c}

\noindent\textbf{Assessment:} Logically valid but extremely slow.

\subsection*{Mathematical goal}
The desired integer is \(\operatorname{lcm}(1,2,\ldots,20)\): the least positive value divisible by every integer from 1 to 20.

\subsection*{How the C code works}
The variable \verb|mult| begins at one. The \verb|for| loop tests \verb|mult%i| for each divisor up to \verb|limit=20|. Whenever a test fails, \verb|mult| increases and the divisor index is reset so that the new candidate is checked again. After all divisibility tests pass, the code prints the candidate.

\subsection*{Correctness and limitations}
The first qualifying candidate is 232792560, but checking essentially every preceding integer is computationally expensive. A much faster implementation would multiply the highest prime powers not exceeding 20. The surrounding \verb|done| loop is not needed to express the search.


\subsection*{Correction to historical C --- Problem 5}
{\color{red}\textbf{Correction:} Unbounded candidate-by-candidate searching is replaced by successive least common multiples computed using the Euclidean greatest common divisor.}
